\documentclass[tikz,border=6pt]{standalone}
% 공통 프리앰블 — PM Day1 교재 그림 (TikZ standalone)
\usepackage{fontspec}
\setmainfont{Apple SD Gothic Neo}
\setsansfont{Apple SD Gothic Neo}
\setmonofont{Menlo}
\usepackage{amssymb}
\usepackage{tikz}
\usetikzlibrary{arrows.meta,positioning,calc,shapes.geometric,fit,backgrounds,matrix,decorations.pathreplacing}
\definecolor{navy}{HTML}{1F3A5F}
\definecolor{teal}{HTML}{2A7F62}
\definecolor{rust}{HTML}{C8553D}
\definecolor{sand}{HTML}{E8E1D5}
\definecolor{ink}{HTML}{222222}
\definecolor{grey}{HTML}{7A7A7A}
\definecolor{lightgrey}{HTML}{EFEFEF}
\definecolor{navylight}{HTML}{DCE6F2}
\definecolor{teallight}{HTML}{DDF0E8}
\definecolor{rustlight}{HTML}{F6E0DA}
\tikzset{
  every node/.style={font=\small, text=ink},
  box/.style={draw=navy, thick, rounded corners=2pt, align=center, inner sep=5pt, fill=white},
  boxfill/.style={box, fill=navylight},
  boxteal/.style={draw=teal, thick, rounded corners=2pt, align=center, inner sep=5pt, fill=teallight},
  boxrust/.style={draw=rust, thick, rounded corners=2pt, align=center, inner sep=5pt, fill=rustlight},
  boxgrey/.style={draw=grey, thick, rounded corners=2pt, align=center, inner sep=5pt, fill=lightgrey},
  arr/.style={-{Stealth[length=2.5mm]}, thick, navy},
  arrteal/.style={-{Stealth[length=2.5mm]}, thick, teal},
  arrrust/.style={-{Stealth[length=2.5mm]}, thick, rust},
  arrgrey/.style={-{Stealth[length=2.5mm]}, thick, grey},
  lbl/.style={font=\footnotesize, text=grey, align=center},
  src/.style={font=\scriptsize, text=grey, align=left},
  title/.style={font=\normalsize\bfseries, text=navy, align=center},
}

\begin{document}
\begin{tikzpicture}
\def\W{150mm}
% 축
\draw[grey, very thick, {Stealth[length=3mm]}-{Stealth[length=3mm]}] (0,0) -- (\W,0);
\node[lbl, anchor=north west, font=\small\bfseries, text=teal] at (0,-2mm) {되돌림 비용 낮음\\→ 빨리 결정하고 틀리면 고친다\\→ 실험 · 반복 (적응형)};
\node[lbl, anchor=north east, font=\small\bfseries, text=rust, align=right] at (\W,-2mm) {되돌림 비용 높음\\→ 결정 전에 최대한 안다\\→ 앞단(front-end) · 동결 (예측형)};
\node[lbl, anchor=north, font=\small\bfseries, text=navy] at (\W/2,-2mm) {하이브리드 — 어느 결정을 언제까지\\되돌릴 수 있게 남기고, 어느 결정을 언제 동결하는가};
% P1
\draw[teal, line width=3pt, opacity=0.5] (5mm,18mm) -- (60mm,18mm);
\node[boxteal, anchor=south west, text width=62mm, font=\footnotesize] at (5mm,21mm) {\textbf{P1 플랫폼 (소프트웨어) — 하이브리드}\\스프린트 1회 $\approx$ 1.41억 (44명 $\times$ 0.5개월 $\times$ 640만 원)\\→ 잘못된 결정 하나의 되돌림 비용 1.4\textasciitilde 2.8억: 실험 가능\\2주 스프린트 36회 · 요구사항 월 1회 동결};
\fill[teal] (60mm,18mm) circle (1.2mm);
% P1 안의 되돌릴 수 없는 두 지점
\node[rust, font=\small] at (108mm,18mm) {$\blacktriangle$};
\node[lbl, text=rust, anchor=south, align=center, font=\scriptsize] at (108mm,21mm) {P1 안의 예외 ①\\G2 데이터 전환 정본 결정\\(2026-11-27; 되돌리면\\610개 테이블 재매핑 3.8억·6주)};
\node[rust, font=\small] at (135mm,18mm) {$\blacktriangle$};
\node[lbl, text=rust, anchor=south, align=center, font=\scriptsize] at (135mm,21mm) {P1 안의 예외 ②\\G3 컷오버\\(34시간, 615개 매장\\펌웨어 일괄배포)};
\draw[teal, dashed] (60mm,18mm) -- (108mm,18mm); \draw[teal, dashed] (108mm,18mm) -- (135mm,18mm);
\node[lbl, font=\scriptsize, anchor=north, text=rust] at (121mm,17mm) {이 두 지점 앞은 예측형으로: 리허설 3회 · 롤백 검증 · point of no return 명시};
% P2
\draw[rust, line width=3pt, opacity=0.5] (95mm,-27mm) -- (\W-5mm,-27mm);
\fill[rust] (128mm,-27mm) circle (1.2mm);
\node[boxrust, anchor=north east, text width=66mm, font=\footnotesize] at (\W-5mm,-30mm) {\textbf{P2 물류 자동화 설비 — 예측형}\\AS/RS 12,400 로케이션 · 셔틀 42대\\설비 발주 확정(T0, 2026-04-24) 후 선급금 29.6억 잠김,\\리드타임 34주(41주 가능성)는 압축되지 않는다\\→ 처리량 4,200 케이스/시간 사양을 발주 \textbf{전}에 확정. 스프린트 무의미};
\node[lbl, anchor=north west, text width=66mm, align=left] at (0,-30mm) {같은 프로그램 안의 두 프로젝트가 축의 반대편에 있다.\\→ 프로그램 하나에 방법론 하나를 강제하면 안 된다.\\[3pt]ONE ONDAM의 다섯 리듬:\\P1 하이브리드 / P2 예측형 /\\P3 파일럿 후 확대(적응형 요소) /\\P4 규제 예측형(기한 고정) /\\P5 운영 이관 예측형(SAC 기준)};
\node[src, anchor=north west] at (0,-64mm) {출처: 교재 2.6절 "되돌림 비용이 방법을 결정한다"의 서술과 사례 문서 수치를 바탕으로 재구성. 축 위 위치는 상대적 개념도.};
\end{tikzpicture}
\end{document}
